% ============================================================================
% APP_ADVANCED_LEARNING.TEX - Lampiran E: Topik Lanjutan & Sumber Belajar
% ============================================================================

\chapter{Topik Lanjutan \& Sumber Belajar}

Untuk mahasiswa yang ingin mendalami pengembangan game lebih jauh, lampiran ini menyediakan panduan mengenai topik-topik tingkat lanjut yang sering digunakan di industri, serta rekomendasi sumber belajar eksternal yang berkualitas.

\section{Advanced Design Patterns}

Design patterns adalah solusi umum untuk masalah yang sering terjadi dalam desain software. Dalam pengembangan game, pola-pola berikut sangat krusial:

\subsection{State Pattern}
Digunakan untuk mengelola state karakter atau game yang kompleks (contoh: Idle, Run, Jump, Attack). Menggantikan `switch` statement yang raksasa dengan class terpisah untuk setiap state.
\begin{itemize}
    \item \textbf{Kegunaan}: AI Controller, Character State Machine.
\end{itemize}

\subsection{Observer Pattern (Events)}
Memungkinkan satu objek untuk memberitahu banyak objek lain bahwa sesuatu telah terjadi tanpa harus "tahu" siapa pendengarnya. Di Unity, ini sering diimplementasikan menggunakan C\# Events atau `ScriptableObject` Events.
\begin{itemize}
    \item \textbf{Kegunaan}: Update UI saat darah berkurang, Achievement system.
\end{itemize}

\subsection{Object Pooling}
Alih-alih membuat (`Instantiate`) dan menghancurkan (`Destroy`) objek terus-menerus (yang membebani CPU dan Memory), kita "meminjam" objek dari kolam (pool) dan mengembalikannya saat tidak terpakai.
\begin{itemize}
    \item \textbf{Kegunaan}: Peluru, musuh yang spawn berulang kali, partikel effect.
\end{itemize}

\subsection{Command Pattern}
Mengubah "permintaan" menjadi objek yang berdiri sendiri. Ini memungkinkan fitur seperti Undo/Redo atau antrian input (input buffering) dalam game fighting.

\section{Arsitektur Game yang Scalable}

\subsection{ScriptableObject Architecture}
Menggunakan `ScriptableObject` bukan hanya untuk menyimpan data, tetapi juga untuk komunikasi antar sistem.
\begin{itemize}
    \item \textbf{Data Containers}: Menyimpan stats musuh atau item agar mudah diedit di Inspector tanpa mengubah script.
    \item \textbf{Event Channels}: Membuat aset ScriptableObject sebagai "jalur komunikasi" di mana satu script bisa memicu event, dan script lain (yang tidak saling kenal) bisa mendengarkannya.
\end{itemize}

\subsection{Dependency Injection}
Teknik untuk mengurangi ketergantungan antar kelas (tight coupling). Framework seperti \textbf{VContainer} atau \textbf{Zenject} populer di Unity untuk proyek skala besar.

\section{Optimasi Performa}

\subsection{Unity Profiler}
Jangan menebak penyebab game lambat. Gunakan Profiler untuk melihat apakah bottleneck ada di CPU, GPU, atau Memory (Garbage Collection).

\subsection{Garbage Collection (GC)}
Hindari alokasi memori berlebih di dalam `Update()`. Contoh: jangan buat `new List<>()` atau string concatenation di setiap frame.

\subsection{DOTS (Data-Oriented Technology Stack)}
Teknologi baru Unity (ECS, Job System, Burst Compiler) untuk performa maksimum dengan memproses data secara paralel dan cache-friendly. Cocok untuk simulasi dengan ribuan unit.

\section{Grafis \& Visual}

\subsection{Render Pipelines (URP \& HDRP)}
\begin{itemize}
    \item \textbf{URP (Universal Render Pipeline)}: Ringan, cocok untuk Mobile dan VR.
    \item \textbf{HDRP (High Definition Render Pipeline)}: Grafis realistis high-end untuk PC/Console.
\end{itemize}

\subsection{Shader Graph}
Membuat shader secara visual (tanpa koding) untuk efek seperti air, hologram, atau distorsi panas.

\section{Rekomendasi Jalur Belajar (Learning Path)}

Berikut adalah kurasi sumber belajar eksternal untuk melengkapi buku ini:

\subsection{YouTube Channels Terbaik}
\begin{itemize}
    \item \textbf{Brackeys}: Tutorial dasar hingga menengah yang sangat jelas (Arsip).
    \item \textbf{Code Monkey}: Fokus pada clean code, design patterns, dan tips praktis profesional.
    \item \textbf{Sebastian Lague}: Eksplorasi mendalam tentang procedural generation dan matematika coding.
    \item \textbf{Jason Weimann}: Arsitektur game dan career advice untuk Unity developer.
    \item \textbf{Tarodev}: Tutorial singkat dan padat tentang fitur-fitur spesifik Unity modern.
\end{itemize}

\subsection{Website \& Artikel}
\begin{itemize}
    \item \textbf{Catlike Coding} (\url{catlikecoding.com}): Tutorial teks mendalam tentang rendering dan hex map.
    \item \textbf{Game Programming Patterns} (\url{gameprogrammingpatterns.com}): Buku online gratis wajib baca tentang design patterns.
    \item \textbf{Unity Learn Premium}: Learning pathway resmi dari Unity.
\end{itemize}

\begin{catatan}
Dunia game development berubah cepat. Selalu cek dokumentasi resmi Unity untuk fitur terbaru.
\end{catatan}
